\documentclass[12pt]{article}
\usepackage[utf8]{inputenc}
\usepackage[T1]{fontenc}
\usepackage{hyperref}
\title{Shahmirzadi Language Resources for K-12 Curriculum Development}
\author{Compiled by Claude Sonnet 5.5 (claude-sonnet-5-5) for Andr\'as Kornai}
\date{}
\begin{document}
\maketitle

\section{Introduction}
Shahmirzadi (ISO 639-3 srz; Glottocode shah1253) is a Caspian-type Indo-European language of the Shahmirzad area, Mehdishahr County, Semnan province, Iran. Habib Borjian's study describes it as sharply different from the other vernaculars of the Semnan area and closely related to Mazandarani. The task description gives about 10,000,000 speakers. That figure is not supported by anything found in search, and it is almost certainly wrong. Joshua Project lists only about 1,100 speakers in Iran, and that figure also differs from other estimates. Shahmirzadi must not be confused with Sangesari (Sangsari), a different Semnan language with its own resources.

\section{Primary Educational Resources}
No primers, textbooks, or school materials in Shahmirzadi were found in search. No evidence was found that the language is taught in schools.

\section{Linguistic and Reference Materials}
\begin{itemize}
\item Borjian, ``The Caspian Language of {\v S}ahmirz\=ad'', Journal of the American Oriental Society: \url{https://lockwoodonlinejournals.com/index.php/jaos/article/view/882} (full text: \url{https://lockwoodonlinejournals.com/index.php/jaos/article/download/882/717/1854}). This is the main linguistic description found.
\item Numeral system of Shahmirzadi (Max Planck Institute for Evolutionary Anthropology, numerals database): \url{https://lingweb.eva.mpg.de/channumerals/Shahmirzadi.htm}
\item Endangered Languages Project entry (it appeared in the Shahmirzadi search, but I did not verify its contents): \url{https://endangeredlanguages.com/node/96185}
\item Sangesari Webonary, for the neighbouring but different language Sangesari. Use it only for comparison: \url{https://www.webonary.org/sangesari/?p=9444}
\end{itemize}

\section{Religious and Community Resources}
Joshua Project's page on Shahmirzadi in Iran gives demographic information: \url{https://joshuaproject.net/index.php/people_groups/18860}. No scripture translation or community organization was found.

\section{Digital and Audio Resources}
No audio, video, or corpus resources were found beyond the numerals page listed above.

\section{Curriculum Development Context}
Very little exists. Any native-language curriculum would first need documentation and community-led orthography work, building on Borjian's description. Instruction in the interim would be in Persian. Speaker numbers must be verified independently.

\end{document}